%% sections/01-introduction.tex
%% Section 1: Introduction (~2 pages)
%% Source notes: v5 LaTeX preamble, Gustavo Notes (April-May 2026),
%% Yao-Zheng-Bai (2015) Chapter 1, Joseph Genzer (2025) Section 1.

\section{Introduction}\label{sec:introduction}

\todo{Open with a one-paragraph hook: random matrix theory began in the 1950s
in nuclear physics and has since become central to high-dimensional statistics
and data science. The fundamental object in multivariate statistics is the
sample covariance matrix, whose eigenvalues describe how variance is
distributed across principal directions. When the number of variables \(p\)
is small compared to the number of samples \(n\), the sample covariance is
well behaved. When \(p\) and \(n\) are comparable, classical intuitions break
down and qualitatively new phenomena appear.}

\todo{Second paragraph: scope. This thesis examines the spectral behavior of
Gaussian sample covariance matrices in two regimes:
(i) the low-dimensional regime where \(p\) is fixed and \(n \to \infty\);
(ii) the high-dimensional regime where \(p/n \to y \in (0, \infty)\).
In the high-dimensional regime, we focus on \emph{Johnstone's spiked model}
\[
  \Sigma = \diag(\alpha_1, \dots, \alpha_m, 1, \dots, 1),
\]
in which the population covariance has finitely many distinguished
eigenvalues (``spikes'') above a flat baseline of ones.}

\todo{Third paragraph: main story. The central phenomenon is the
\emph{BBP phase transition}: a spike \(\alpha > 1 + \sqrt{y}\) creates an
isolated outlier in the sample spectrum with informative limiting
eigenvector, while a sub-threshold spike is absorbed into the bulk and its
eigenvector becomes uninformative. This contrast underpins when methods like
principal components analysis and spectral clustering succeed in high
dimensions.}

\todo{Fourth paragraph: roadmap. Section~\ref{sec:foundations} covers the
probability foundations needed for the rest of the thesis. Section~\ref{sec:scm}
develops sample covariance matrices, the Wishart distribution, the
low-dimensional Gaussian and GOE fluctuation regimes, and the
high-dimensional Marchenko-Pastur and Tracy-Widom laws.
Section~\ref{sec:spiked} introduces the spiked covariance model and proves
the BBP phase transition for eigenvalues, discusses fluctuations of spiked
eigenvalues, derives the eigenvector alignment formula, and sketches the
connection to spectral clustering. Section~\ref{sec:conclusion} concludes
with a summary and open questions. Three appendices collect auxiliary
results, simulation methodology, and selected code.}

\clearpage
